It is clear that there is a problem with computational models right now and while there are solutions proposed they still did not have enough traction to truly change this scientific field. This is where \greensloth{} comes into play. The main aspect of \greensloth{} is creating a bridge between scientists that create models and those that will actually try to use them for experimental work, as these two groups often do not coincide. To create this bridge, \greensloth{} will encompass many aspects of computational work that are already being tackled on as mentioned prior but with one key advantage. It will minimise its scope onto kinetic models of photosynthesis.

While models come in many different shapes and sizes, all showing different things, We believe it is important to first reduce the problem at hand to a manageable feat, by only encompassing a single, yet incredibly important, biochemical process. The importance of this process has also not escaped the eyes of the modelling community, as seen by the incredible amount of publication every year~\figref[]{fig:photosynthesisModelQuery}. Sifting through these amounts of publications is hard for both the modellers, that may want to expand their model with other interpretations, and the experimentalists that may want to build upon their work by using the new, popular and useful method of computational modelling.

The \greensloth{} project can be seperated into four main branches \figref{fig:greenSlothPlan}. These branches are Model Creation Pipeline, Model Comparison \& Validation, Database, and Website. Each branch has its own set of tasks and goals that are integral to \greensloth{} to create a comprehensive platform for kinetic models of photosynthesis. All these branches will be explained in the following sections.

\begin{figure*}[!htb]
    \centering
    \includegraphics[width=0.8\textheight, angle=90]{Figures/GreenSloth_Plan.pdf}
    \captionprof{The tree of \greensloth{}.}{A schematic representation of the different aspects of the \greensloth{} project and how they relate to each other. The four main branches are Model Creation Pipeline, Model Comparison \& Validation, Database, and Website. Each branch has its own set of tasks and goals that will be tackled to create a comprehensive platform for kinetic models of photosynthesis.}
    \label{fig:greenSlothPlan}
\end{figure*}

\subsection{The Promise}

As shown with the large amount of publications on kinetic models of photosynthesis every year~\figref[]{fig:photosynthesisModelQuery}, and a large amount that cannot be directly reproduced (see Tiwari et al.~\cite{tiwariReproducibilitySystemsBiology2021}), model creation is not a simple and straightforward task. You start by having a question on real systems that you specify into a problem. From there, you gather knowledge on your system to design a conceptual model. With this conceptual model, you can formalise a mathematical model, where you sometimes have to make assumptions based on your knowledge of the system. You can now either calculate analytical solutions, or more commonly program your model into a computational format to simulate it. Either way, you arrive into results, which you have to validate with your initial or other real systems. If the results do not depict your system well enough, you reevaluate and create a new problem to start the process again. If the results are satisfactory, you have solved your real questions with a mathematical model \figref{fig:modelingScheme}.

While \greensloth{}'s intention is solely based on a model that has already reached the solved stage, it still encapsulates several aspects of the model creation process. The design of a conceptual model, the formulation into a mathematical model, the programming into a computational format, examples of model validation are all part of the \greensloth{} project. Having these aspects semi-standardised, will ensure the quality of models which will help the modellers with the publishing of their models, and keeping them under the \gls{cure} guidelines. Furthermore, it will also help users of models, understand the process of mathematical modelling better, and therefore be able to use models more effectively.

To help the users of models more, \greensloth{} is the central hub to find photosynthesis kin etics models. It provides a way to compare models, see their capabilities and help the user with deciding which model is best suited for their needs. This shoudl eliviate the problem of finding the right model, and opens up the field of mathematical models of photosynthesis to a wider audience. Therefore, \greensloth{} represents the bridge between a model author and a model user, helping both sides with their respective problems.

\begin{figure*}[!htb]
    \centering
    \includegraphics[width=0.7\textwidth]{Figures/modelling.pdf}
    \captionprof{A simple scheme of mathematical modelling.}{A schematic representation of the mathematical modelling process, starting from a real system with a question, to finding a solution using a model. After specifying the problem, knowledge is gathered to design a conceptual model. This should then be formalised into a mathematical mode, by adding assumptions if needed. From there, analytical or computational solutions can be found to arrive at results. These results need to be validated with real systems, and if they are not satisfactory, the process starts again from the problem definition. If they are, the real question has been solved using a mathematical model.}
    \label{fig:modelingScheme}
\end{figure*}

\subsection{Model Creation Pipeline}

To ensure a standardised model format, a pipeline on how to create a model will be created. This pipeline will emcompass a guide to formatting the code itself, automatically export the model infromation, for example the \gls{ode} sytem, parameters, etc., into tables and math expressions, force the user to provide a summary and a recreation of the publication figures with code. All this information will then be stored in one directory, with a Markdown file as its front page. This directory can then be easily uploaded to the \greensloth{} database. This pipeline will be provided as a Python package under the name \greensloth[GreenSlothUtils].

\subsubsection{Model Code}

To ensure a common format for the model code itself, a set requirement for it, will be using the MxlPy package~\cite{vanaalstMxlPyPythonPackage2025}. This package enables a building block approach to building \gls{ode} models, however by keeping the familiar linear flow of Python. Furthermore, it allows for easy simulation, plotting, parameter fitting, and many more functionalities that are crucial while using a model. Because of that, and continuing support from the authors, it is the perfect candidate to be the base of \greensloth{} models.

On top of that, the actual formatting of the code will be standardised to a template that will be provided in the \greensloth{} model creation pipeline. This template will ensure that all necessary information is provided, and that the code is easy to read and understand. It will be continously updated to reflect the best practices in coding and modelling, for example by having the same names for common biological variables, such as ATP, NADPH, etc.

\subsubsection{Glossaries}

To ensure that all models are properly annotated, a set of glossaries will be created. These glossaries will be created for the reactions, the variables, the parameters, and the derived versions of the latter two. These glossaries will provide a standardised way of naming and describing these entities, ensuring that they are easily understandable and comparable between models. Furthermore, they will provide a way to link these entities to external databases, such as \gls{kegg}~\cite{kanehisaKEGGKyotoEncyclopedia2000}. To facilitate this process, the \greensloth[GreenSlothUtils] package will provide functions to automatically extract the information from the MxlPy model. Furthermore, it will alert the user if all the information that the model entails can be found in the glossaries.

These glossaries will come in two forms. One tied with the model itself, as many parameters and variables may be model-specific, due to assumptions made during the model creation. These glossaries will follow how the publication of the model presents the existing information. The second form will be a global glossary, that will provide a standardised way of naming and describing common entities between models. This global glossary will be continuously updated as more models are added to the \greensloth{} database.

\subsubsection{Gather Information}

To make the model information as easily accessible as possible, the \greensloth[GreenSlothUtils] package will provide functions to automatically create a markdown file that will summarise the model. It will automatically create tables from the glossaries, and let the user be able to automatically generate the \gls{ode} system and other needed mathematical expression of the model in \LaTeX{} format. Furthermore, it will provide space for the model validation to remind the user to include them.

This markdown file is a prerequisite for any model in the \greensloth[] database, and shall facilitate the process of publishing the models to the \gls{cure} standards. 

\subsubsection{Packaging it up}

All the information gathered from the previous steps will then be packaged into a single directory, that will contain the model code, the glossaries, the markdown file, and all the code used for the model validation. This directory will then be easily uploadable to the \greensloth{} database and also directly be transformed into a Python package that can be installed and used by anyone. This enables easy sharing and reusing of models, as well as ensuring that all the necessary information is provided.

\subsection{Model Comparison \& Validation}

With the information gathered from the model creation pipeline, the first step to comparing models is already done. The standardised format of the models, as well as the tables and mathematical expressions of critical model information already show what different models were made to show. However, to truly compare models, they need to be validated the same way. Therefore, a set of standardised validation protocols will be created, that will test each model to common practices in \gls{ode} models and photosynthesis work. These protocols will be provided as part of the \greensloth[GreenSlothUtils] package, and will be applicable to any model in the \greensloth[] database.

The first step in validation is to ensure that the model is mathematically sound. This includes checking for mass balance, energy balance, and other common checks that are applicable to \gls{ode} models. One crucial aspect, for example, is to ensure that the model does not produce negative concentrations, as this is not physically possible. This step can colloquially be called "Sanity Checks", as it ensures that the model behaves in a realistic and sane manner.

The second step in validation, is the recreation of the publication figures. This step is crucial to ensure that the model behaves as expected and presented, and that the code provided is correct. This means that the code used to create these figures needs to be manually curated to ensure true results. With the recreation of these figures, a score can be set to see how well the model performs compared to the publication. Depending on the model curated and the information provided, help from the authors may be needed to ensure the best possible recreation of the figures. However,if additional information is needed to recreate the figures in the publication, this is proof that the model does not follow the \gls{cure} principles, and therefore the version on \greensloth[] will entail the revisions needed and give the community the best possible version of the model.

The third and final step in validation, is to create a set of standardised tests that are applicable on a wide variety of photosynthesis models. These tests will be based on expert knowledge, from modellers to experimentalists, and will show the extensive simulation possibilities of a model. In this step, a model may not be able to pass all the tests, as it may not have been designed to do so. Therefore, a rating on how "good" a model is, cannot be done using this step as it can with the former mentioned. However, this will give a good overview of the model capabilities, and will allow for easier comparison between models and should help users determine which model is best suited for their needs. One example for these tests include "injecting" the simple \gls{fvcb} model~\cite{voncaemmererSteadystateModelsPhotosynthesis2013} into the model and see how it simulates the commonly used relation of \gls{A} and \gls{Ci}.

\subsection{Database}

Having a centralised database to store and share models is crucial to ensure that they are easily accessible and reusable. This database will be the backbone of \greensloth{}, as it will provide a way to store the models. In the beginning, this database will be a simple file-based system, where each model is stored in a separate directory. This can easily be hosted on GitHub, which provides version control and collaboration features. Furthermore, it is a well-known and used platform that strives on an open-source nature. Exactly that, what \greensloth{} is about. This openness, also allows access to all the models inside the database, for users, but also other parts of \greensloth{}, like the website \cleverref[~\nameref*{sec:website}]{Section}{sec:website}. However, as the number of models grows, this system may become unwieldy. Therefore, a more robust database system will be implemented, that will still let the website access all models easily.

\subsection{Website} \label{sec:website}

Sharing models is no easy task. Only using publications can be tedious, as the code may not be available, or the model is not properly explained. Therefore, a centralised platform to collect and present models, would be beneficial to the scientific community. We believe that the easiest way to create such platform is in a web-based format, where anyone can access it from anywhere. This platform will present models in a standardised format, that includes a summary, a table of model information, and the recreation of the publication figures. Furthermore, having a central location with all the models allows for easier comparison between them. 

Web development and design has become an integral part of modern communication, boasting around 74\% of the global population being internet users in 2025~\cite{MeasuringDigitalDevelopment}. From social media to spreading knowledge, from entertainment to education, and many more, websites shape how people interact with their politics, their art, their community, and themselves~\cite{DataReportalGlobalDigital}. Therefore, a great deal of research has been done to optimize and cater web pages to the public. Good we design can be separated into four categories~\cite{nawirImpactWebsiteUsability2024}: ease of navigation, site speed, clarity and conciseness of content, and visually appealing design. There are many resources that help web developers deal with these aspects, such as books and guidelines~\cite{beairdPrinciplesBeautifulWeb2020}, further updates to existing basic web development languages (HTML, CSS, JavaScript), additional web development frameworks (React, Angular, Vue.js) that help create better web design, and many more. Web development is as popular as never before.

Therefore, giving access to \greensloth{} through a well-designed website is crucial to its success. This rise to popularity of web development can be seen in science, as more and more scientific databases and tools are being made available through web interfaces. A few examples are \gls{blast}~\cite{altschulBasicLocalAlignment1990}, \gls{kegg}~\cite{kanehisaKEGGKyotoEncyclopedia2000}, \gls{bigg}~\cite{kingBiGGModelsPlatform2016}, or the aforementioned BioModels, with many more examples in various scientific fields. This rise in popularity also shows the advantage of web interfaces over other platforms, such as applications or software packages. They are much more accessible and easier to use for the public, as they do not require any installation or specific operating system. Furthermore, they advocate an open-source approach, and therefore align with the \gls{fair} and \gls{cure} principles, making scientific data and tools more accessible and reusable for everyone.